% !TEX root = ../main.tex

\begin{acknowledgements}
在中山大学计算机学院攻读硕士学位的三年，对我来说毫无疑问是一段格外珍贵的旅程。作为一名非全日制学生，我需要在工作与学业之间寻找平衡，整个学业的完成，包括本篇论文，离不开许多人的理解、支持和帮助，要感谢的人和事很多，简要概括如下。

一谢师恩。首先，我要衷心感谢我的导师。由于我是非全日制学生，无法像全日制同学那样全天候待在实验室，但老师从未因此降低对我的要求，反而给予了更多的包容与耐心。从选题方向的确定，到研究方法的指导，再到论文的逐字修改，老师始终以严谨的治学态度和深厚的学术功底引领我前行。每当我因工作繁忙而进度滞后时，老师总能给予充分的理解和及时的督促。这份宽容与信任，是我能够坚持完成学业的最大动力和助力。

二谢同窗。感谢实验室的各位小伙伴，同时也很荣幸能够结识到学术上严谨，思维上活跃，生活上有趣的优秀的你们。虽然我在校时间有限，但每一次组会、每一次小组分享，大家都毫无保留地交流最新的前沿技术动态和研究心得。正是在这种开放、互助的氛围中，让我在有限的在校时间里，能够既接触到Transformer、Diffusion等模型的前沿研究，也能接触到NVIDIA的前沿工业实践，拓宽了我的技术视野。在此特别感谢几位师兄师姐在我访问实验室GPU资源受阻时给予了指导，在我研究方向出现漏洞时及时提醒，那些深夜的头脑风暴让我受益匪浅。

三谢良助。我还要感谢学院教务和行政岗位的老师们。作为一名非全日制学生，我在选课、注册、答辩申请等各个环节都遇到过不少因工作与学校作息冲突带来的不便。每一次打电话或发邮件咨询，老师们都耐心细致地为我解答，帮我协调时间、解决问题。你们的敬业与友善，让我在奔波于公司与校园之间的日子里，感受到了学校的温暖。

再谢授业。感谢研究生期间每一位授课老师，感谢你们在课堂上的精彩讲授。这些宝贵的课堂知识，不仅帮助我将工作中的实践经验与理论知识有机结合，更拓展了我的技术视野，深化了我对计算机学科的理解。每一门课程的学习，都为我后来的研究工作打下了坚实的基础，为本篇论文的完成添砖加瓦。我相信，它们也必将给我之后的工作学习带来意想不到的助力。

终谢母校。感谢中山大学提供的宝贵平台与便利资源。三年的时光碎片里，有超算中心实验室里热烈的讨论声，有第二教学楼窗外树影的斑驳交错，有松涛园手撕鸡的咸香，还有逸仙路草坪雨后泥土的清新气息。身处中大校园，无论是安静的自习场所、设备齐全的实验室，还是丰富的电子图书资源和学术数据库，都为我的学习和研究提供了极大的支持。能在康乐园度过这段求学时光，是我的荣幸。

三年的时光虽短，但在中大收获的知识、眼界与情谊，将伴随我走得更远。

\end{acknowledgements}
